\documentclass[12pt,a4paper]{article}
\usepackage[T1]{fontenc}
\usepackage[utf8]{inputenc}
\usepackage[brazil,english]{babel}
\usepackage{amsmath,amssymb,amsfonts,amsthm}
\usepackage{geometry}
\usepackage{graphicx}
\usepackage{booktabs}
\usepackage{array}
\usepackage{fancyhdr}
\usepackage{titlesec}
\usepackage{xcolor}
\usepackage{setspace}
\usepackage{mathtools}
\usepackage{bm}
\usepackage{hyperref}
\usepackage{enumitem}
\usepackage{tcolorbox}
\usepackage{mdframed}

\geometry{a4paper, top=3cm, bottom=2.5cm, left=3cm, right=2cm, headheight=25pt}
\setstretch{1.35}

\definecolor{ufamblue}{RGB}{0,51,102}
\definecolor{ufamgreen}{RGB}{0,102,51}
\definecolor{resultbg}{RGB}{240,248,255}

\hypersetup{colorlinks=true, linkcolor=ufamblue, urlcolor=ufamblue, citecolor=ufamblue}

\pagestyle{fancy}
\fancyhf{}
\fancyhead[L]{\small\textcolor{ufamblue}{\textbf{PPGEE/UFAM} --- PGENE601}}
\fancyhead[R]{\small\textcolor{ufamblue}{Resumo cr\'itico --- Survey ECPS}}
\fancyfoot[C]{\small\thepage}
\renewcommand{\headrulewidth}{0.6pt}
\renewcommand{\headrule}{\hbox to\headwidth{\color{ufamblue}\leaders\hrule height \headrulewidth\hfill}}

\titleformat{\section}[block]
  {\large\bfseries\color{ufamblue}}
  {\thesection.}{0.8em}{}
  [\vspace{-0.3em}\textcolor{ufamblue}{\rule{\linewidth}{0.6pt}}]
\titlespacing*{\section}{0pt}{2.5ex}{1.5ex}

\titleformat{\subsection}{\normalsize\bfseries\color{ufamgreen}}{\thesubsection}{0.8em}{}
\titlespacing*{\subsection}{0pt}{1.5ex}{0.8ex}

\tcbset{
  infobox/.style={
    colback=resultbg, colframe=ufamblue, fonttitle=\bfseries\small,
    title=Reference, boxrule=0.8pt, arc=3pt,
    left=6pt, right=6pt, top=4pt, bottom=4pt
  }
}

\begin{document}
\selectlanguage{english}

%==============================================================
% CAPA
%==============================================================
\begin{titlepage}
\centering
\vspace*{0.5cm}

{\color{ufamblue}\rule{\linewidth}{2pt}}\\[0.4cm]
{\Large\bfseries\color{ufamblue} UNIVERSIDADE FEDERAL DO AMAZONAS}\\[0.2cm]
{\large\color{ufamblue} Programa de P\'os-Gradua\c{c}\~ao em Engenharia El\'etrica --- PPGEE}\\[0.1cm]
{\large\color{ufamblue} Linha 1 --- Sistemas Inteligentes e Microeletr\^onica}\\[0.4cm]
{\color{ufamblue}\rule{\linewidth}{2pt}}\\[2.5cm]

{\LARGE\bfseries Critical Summary}\\[0.4cm]
{\Large\bfseries ``A Survey on Automated Symbolic Verification and its}\\[0.1cm]
{\Large\bfseries Application for Synthesizing Cyber-Physical Systems''}\\[0.5cm]
{\large L.\,C.\ Cordeiro, E.\,B.\ de Lima Filho, I.\,V.\ Bessa}\\[0.1cm]
{\normalsize IET Cyber-Physical Systems: Theory \& Applications, 2020}\\[3cm]

\begin{tabular}{rl}
  \textbf{Disciplina:} & PGENE601 --- Verifica\c{c}\~ao e S\'intese Autom\'atica de \\
                        & Sist.\ Ciber-F\'isicos e Embarcados \\[0.3cm]
  \textbf{Docente:}    & Prof.\ Dr.\ Lucas C.\ Cordeiro \\[0.3cm]
  \textbf{Discente:}   & Matheus Serr\~ao Uch\^oa \\[0.3cm]
  \textbf{N\'ivel:}    & P\'os-Gradua\c{c}\~ao --- Mestrado \\[0.3cm]
  \textbf{Data:}       & Manaus, \today \\
\end{tabular}

\vfill
{\color{ufamblue}\rule{\linewidth}{1.5pt}}
\end{titlepage}

%==============================================================
% RESUMO
%==============================================================

\begin{tcolorbox}[infobox]
\textbf{Full reference.} L.~C.~Cordeiro, E.~B.~de Lima Filho, I.~V.~Bessa. ``A Survey on
Automated Symbolic Verification and its Application for Synthesizing Cyber-Physical Systems.''
\emph{IET Cyber-Physical Systems: Theory \& Applications}, vol.~5, no.~1, pp.~1--24, 2020.
DOI: 10.1049/iet-cps.2018.5006.
\end{tcolorbox}

\section{Objectives}

The survey addresses the growing reliability demands placed on \emph{embedded and
cyber-physical systems} (ECPS) --- systems that combine discrete computation with continuous
physical dynamics and are now pervasive in automotive, medical, industrial, and IoT
applications. The authors set three explicit goals: (i)~briefly describe the state of the art
in ECPS modeling, verification, and synthesis; (ii)~provide a technical overview of
\emph{symbolic} verification methods, with particular emphasis on bounded model checking (BMC)
and its extensions, which the authors identify as their own ``vision for future ECPS
verification and synthesis schemes''; and (iii)~present open challenges in ECPS verification
and synthesis and clarify current trends, which the authors argue are largely driven by ECPS
growth, the Internet of Things, and machine learning (ML). The survey's scope is explicitly
restricted to \emph{automated symbolic} verification and synthesis; simulation-based and purely
statistical approaches are mentioned only for contrast, not reviewed in depth.

\section{Methods}

Two distinct notions of ``method'' operate in this paper: the survey methodology itself, and
the technical verification/synthesis methods it surveys.

\subsection{Survey methodology}

The authors conducted a systematic literature search across three digital libraries --- Google
Scholar, Web of Science, and Scopus --- following a three-stage protocol of \emph{searching},
\emph{screening}, and \emph{synthesis}. To mitigate bias inherent to automated database
queries, search results were cross-validated against a manually curated ``control list'' of
papers from peer-reviewed venues already known, \emph{a priori}, to be relevant to
symbolic verification and synthesis for ECPS. Screening was guided by inclusion criteria
centered on the core topic keywords (e.g., ``Internet-of-Things'', ``Concurrent Systems'',
``Hybrid Systems'', ``Control Systems''). The authors candidly acknowledge, in a dedicated
limitations section, that this screening stage remains partly subjective and that their search
may not cover the entire ECPS literature.

\subsection{Technical methods surveyed}

ECPS are first formalized as \emph{hybrid automata}: a tuple $S=(X,X_0,U,Y,r,v)$ of reachable
states, initial states, inputs, outputs, a transition relation, and an output map. This
modeling layer is illustrated throughout the paper with a running example, the \emph{Sim3Tanks}
three-tank system, whose safety and liveness requirements are expressed both in computation
tree logic (CTL) and in linear-time temporal logic (LTL). Building on this foundation, the
survey reviews six families of symbolic \textbf{verification} technique:

\begin{enumerate}[leftmargin=1.8cm]
  \item \textbf{Bounded Model Checking (BMC).} Unrolls a transition system to a fixed depth $k$
  and checks satisfiability of the resulting verification condition via SAT or SMT solvers.
  Sound, but --- absent a known completeness threshold --- inherently incomplete.
  \item \textbf{$k$-induction.} Strengthens BMC into a complete proof technique by combining a
  base case with an inductive step over automatically inferred loop invariants, removing the
  need for users to manually supply them.
  \item \textbf{Property-Directed Reachability (IC3/PDR).} An alternative inductive-invariant
  search that is guided by counterexamples-to-induction (CTIs); the survey reports that IC3
  scales better than $k$-induction on certain benchmark classes.
  \item \textbf{Craig interpolation.} Derives an over-approximation (an interpolant) of
  reachable states directly from an SMT solver's unsatisfiability proof, offering a route to
  invariant generation that complements $k$-induction and IC3.
  \item \textbf{Abstraction.} Predicate abstraction with counterexample-guided refinement
  (CEGAR) and abstract interpretation reduce model size by tracking only relevant predicates or
  numerical intervals, trading precision for scalability (representative tools: SLAM, SATABS,
  Astr\'ee, PolySpace).
  \item \textbf{Path-based symbolic execution.} Explores individual program paths (e.g., KLEE),
  useful for coverage-driven test generation but subject to path explosion on large programs.
\end{enumerate}

For \textbf{synthesis}, the survey centers on \emph{Counter-Example Guided Inductive Synthesis}
(CEGIS): an iterative \textsc{Synthesize}--\textsc{Verify} loop that refines a candidate
controller against a correctness specification using counterexamples as new constraints, until
a verified solution is found or the search fails. Its application to ECPS is illustrated through
\textbf{DSVerifier}/\textbf{DSSynth}, an ESBMC-based front-end that instruments digital-controller
code (direct, delta, and cascade realizations) with finite-word-length semantics and checks
overflow, limit-cycle, timing, stability, and minimum-phase properties --- the same tool
examined earlier in this course's exercise list on DSVerifier and $k$-induction, which gives
this summary a direct, concrete anchor to the practical material already covered.

\section{Results (Key Findings)}

As a survey, this paper's ``results'' are synthesized findings rather than experimental data:

\begin{itemize}[leftmargin=1.8cm]
  \item \textbf{Eleven concrete challenges} for embedded/ECPS verification and synthesis are
  catalogued: time and energy constraints; handling concurrent software; platform restrictions;
  security issues; verifying code modules embedded in larger structures; legacy designs;
  support for multiple programming languages and interfaces; correct code instrumentation;
  non-linear/non-convex optimization; incorporation of new checks as systems evolve; and
  provision for system evolution over successive versions.
  \item These challenges are distilled into \textbf{seven open Research Problems (RP1--RP7)}:
  suitable SMT encodings for realistic ECPS background theories; scaling BMC-based verification
  of multi-threaded software; proving both correctness \emph{and} timeliness via stronger
  $k$-induction/invariant techniques; verifying system-level (not merely code-level) properties;
  broadening tool, language, and application support; more robust automated formal synthesis;
  and simultaneous verification of low- and high-level properties within a single framework.
  \item A cross-reference table explicitly links each research problem to the challenges it
  addresses, showing that the most demanding problems (RP1 and RP7) cut across nearly every
  challenge category identified.
  \item A dedicated section looks \textbf{beyond code correctness}, discussing SMT-based
  non-convex optimization (via a CEGIS variant, CEGIO) and introducing \emph{antifragility}:
  systems designed not merely to resist faults but to strengthen from exposure to them --- an
  idea the authors connect to a future combination of ML with symbolic verification loops.
  \item The authors explicitly report their own \textbf{limitations}: bias inherent to
  automated digital-library search, a screening stage acknowledged as partly subjective, and a
  scope restricted to symbolic (as opposed to statistical or purely simulation-based) methods.
\end{itemize}

\section{Conclusion}

The survey concludes that, while BMC, $k$-induction, IC3, Craig interpolation, abstraction, and
path-based symbolic execution are individually mature techniques, there remains limited
published evidence of their systematic \emph{combination} (e.g., $k$-induction strengthened
with learned invariants, or IC3 combined with induction) applied to ECPS specifically. Formal
synthesis via CEGIS, in turn, is shown to be fundamentally dependent on --- and bottlenecked by
--- the soundness of the underlying verification models it queries for counterexamples. The
authors' own future-work agenda mirrors this gap directly: a \emph{unified} framework that
jointly addresses code verification, behavior correctness, environment adaptation, and system
evolution, with machine learning incorporated to guide search and strengthen invariants. This
positions the survey less as a finished map of the field and more as an explicit research
agenda --- one that the tools examined earlier in this course (ESBMC, DSVerifier, and the
$k$-induction algorithm) already partially instantiate, but that the authors argue is still far
from the ``unified, high- and low-level, correct-by-construction'' vision they set out for ECPS.

\end{document}
